\Article{22}
Die vorangehenden Versuche wurden vor allem unternommen, um das Gesetz der magnetischen Wirkung über jeden Zweifel zu erheben. Außerdem sollten sie zeigen, wie viele Terme der Reihe man berücksichtigen sollte und welche Genauigkeit die Versuche zulassen. Sie haben gezeigt, dass, wenn wir bei den Abständen das Vierfache der Nadellänge nicht \SourceLeaf{115}unterschreiten, zwei Terme der Reihe genügen\HistoricalNote{12}. Die Unterschiede, die sich aus der Rechnung ergeben haben, dürfen übrigens keineswegs bloß als Beobachtungsfehler gelten. Mehrere Vorkehrungen, von deren Anwendung wir eine noch bessere Übereinstimmung erhoffen dürfen, waren damals noch nicht vorbereitet. Dazu gehören Korrekturen wegen der stündlichen Schwankungen der Intensität des Erdmagnetismus. Man muss sie mithilfe einer weiteren Nadel zum Vergleich berücksichtigen, nach dem Vorbild des Verfahrens, das wir in Artikel $10$, II. besprochen haben. Um den Wert des Erdmagnetismus zu bestimmen, soweit er sich aus \textit{diesen} Versuchen ableiten lässt, fügen wir eine Übersicht der übrigen einschlägigen Versuche hinzu.

Der Wert des Bruchs $\frac{\theta}{Tm}$ für die erste Nadel und den Faden, an dem sie hing, wurde nach dem in Artikel $8$ angegebenen Verfahren zu $=\frac{1}{251{,}96}$ bestimmt. Daraus folgt
\[
\frac MT=0{,}0436074
\]
Dieser Zahl liegt der Meter als Einheit der Abstände zugrunde. Ziehen wir als Einheit den Millimeter vor, so müssen wir diese Zahl mit dem Kubus von tausend multiplizieren. Das ergibt
\[
\frac MT=43607400
\]

Für die zweite Nadel ergaben die am $28$. Juni durchgeführten Versuche, die denen für eine andere Nadel in Artikel $11$ ganz entsprachen, das folgende Ergebnis, wenn man den Millimeter, das Milligramm und die Sekunde der mittleren Sonnenzeit als Einheiten nahm.
\[
TM=135457900
\]
Eliminiert man daraus die Größe $M$, so folgt
\[
T=1{,}7625
\]
